\pdfminorversion=4
\documentclass[]{article}
\usepackage[utf8]{inputenc}
\usepackage{amssymb,latexsym,amsmath}
\usepackage[a5paper,top=2cm,bottom=2cm,left=2cm,right=2cm,marginparwidth=1.75cm]{geometry}
\usepackage{graphicx}
\usepackage[colorlinks=true, allcolors=blue]{hyperref}
\usepackage{helvet}
\usepackage{tikz}
\usepackage{calc}
\renewcommand{\familydefault}{\sfdefault}
\newcommand{\ligature}[2]{%
  \begingroup
  \leavevmode
  \setbox0=\hbox{#1}%
  \setbox1=\hbox{#2}%
  \tikz[baseline=(base.base)]{
    \node[inner sep=0pt] (base) {\strut #1~#2};
    \coordinate (start) at ([xshift=\wd0-0.1em,yshift=0.5ex]base.south west);
    \coordinate (end) at ([xshift=\wd0+0.5em,yshift=0.5ex]base.south west);
    \draw[line width=0.25pt]
      (start) .. controls +(0.25em,-0.15em) and +(-0.25em,-0.15em) .. (end);
  }%
  \endgroup
}
\newcommand{\ligaturethree}[3]{%
  \begingroup
  \leavevmode
  \setbox0=\hbox{#1}%
  \setbox1=\hbox{#2}%
  \setbox2=\hbox{#3}%
  \tikz[baseline=(base.base)]{
    \node[inner sep=0pt] (base) {\strut #1~#2~#3};
    \coordinate (start1) at ([xshift=\wd0-0.1em,yshift=0.5ex]base.south west);
    \coordinate (end1) at ([xshift=\wd0+0.5em,yshift=0.5ex]base.south west);
    \draw[line width=0.25pt]
      (start1) .. controls +(0.25em,-0.15em) and +(-0.25em,-0.15em) .. (end1);
    
    \coordinate (start2) at ([xshift=\wd0+\wd1+0.2em,yshift=0.5ex]base.south west);
    \coordinate (end2) at ([xshift=\wd0+\wd1+0.8em,yshift=0.5ex]base.south west);
    \draw[line width=0.25pt]
      (start2) .. controls +(0.25em,-0.15em) and +(-0.25em,-0.15em) .. (end2);
  }%
  \endgroup
}


\begin{document}

\hypertarget{h001}{}
\section*{001. Sur la croix}
\addcontentsline{toc}{section}{001. Sur la croix}
\label{h:001}
\index{Fondement de l'Église}

\vspace{-0.8em}
\noindent{\small\textit{Rm 5:8 - Ep 2:8-9 - 2Co 5:21 - Ep 5:25 - Jn 14:2-3}}

\vspace{0.8em}
\medskip
\noindent\textbf{1.} 1. La mort sur la croix il n'existe pas egal,\\
Là, Dieu révéla son pardon.\\
Parfait et sans faute, Il fut crucifié,\\
Pour sauver l'homme loin de la foi.

\medskip
\indent\textbf{Quel amour, grand amour,}\\
\indent\textbf{\ligaturethree{Qu'on}{trouve}{en} Jésus seul, à jamais!}

\medskip
\noindent\textbf{2.} L'amour sur la croix, seul Christ a pu l'accomplir,\\
Personne d'autre ne pouvait pas.\\
En Lui, la justice de Dieu fut donnée,\\
Et la paix avec Dieu retrouvée.

\medskip
\noindent\textbf{3.} Le don sur la croix a payé tout nos péchés,\\
Son oeuvre nous a délivrés.\\
Par grâce eternel, Il nous a rachetés,\\
Et au ciel, Il va nous amener.

\end{document}
